\uuid{1iCr}
\titre{Limite}
\auteur{bodin ; hanani ; mzari}
\organisation{exo7}
\licence{CC BY-NC-SA 4.0}
\datecreate{2020-01-20}
\niveau{L1}
\module{Analyse}
\chapitre{Suite}
\sousChapitre{Autre}
\difficulte{2}
\theme{Suites, limite}
\exo7qcmid{lille-1.suites.018}

\contenu{
\begin{question}
\qclassification{121.00}
\qtags{motcle=limite}
Soit $(u_n)$ une suite réelle. On suppose que $\displaystyle |u_{n+1}-1|\leq \frac{1}{2}|u_n-1|$ pour tout $n\geq 0$. Que peut-on en déduire ?
\begin{answers}
    \bad{La suite $(u_n)$ est convergente et $\displaystyle \lim _{n\to +\infty}u_n=0$.}
    \bad{La suite $(u_n)$ est divergente.}
    \good{Pour tout $n\geq 1$, $\displaystyle |u_n-1|\leq \frac{1}{2^n}|u_0-1|$.}
    \good{$\displaystyle \lim _{n\to +\infty}u_n=1$}
\end{answers}
\begin{explanations}
On vérifie par récurrence que $\displaystyle |u_n-1|\leq \frac{1}{2^n}|u_0-1|$, et donc, par passage à la limite, $\displaystyle \lim _{n\to +\infty}(u_n-1)=0$. C'est-à-dire, $\displaystyle \lim _{n\to +\infty}u_n=1$.
\end{explanations}
\end{question}
}
